% Source: https://artofproblemsolving.com/community/c3h1567867_mathcounts_cdr_chapter
% Stable identity: aops-topic-1567867-q46
We need integers with $x^2+y^2=169=13^2$, i.e.\ $y^2=169-x^2=(13-x)(13+x)$, with $|x|\le13$. Checking $|x|=0,1,\dots,13$, the value $169-x^2$ is
\[
169,168,165,160,153,144,133,120,105,88,69,48,25,0,
\]
and it is a perfect square only for $|x|=0,5,12,13$ (values $169,144,25,0$). So the solutions are
\[
(0,\pm13),\ (\pm13,0),\ (\pm5,\pm12),\ (\pm12,\pm5),
\]
which is $2+2+4+4=\boxed{12}$ points. They come from the Pythagorean triple $5,12,13$ and the trivial ones on the axes.

\begin{asy}
usepackage("amsmath");
size(5cm);
defaultpen(fontsize(10pt));
draw(circle((0,0),13),linewidth(0.9));
draw((-15,0)--(15.5,0),gray,Arrow(TeXHead));draw((0,-15)--(0,15.5),gray,Arrow(TeXHead));
pair[] L={(13,0),(-13,0),(0,13),(0,-13),(5,12),(-5,12),(5,-12),(-5,-12),(12,5),(-12,5),(12,-5),(-12,-5)};
for(pair p:L) dot(p,red+linewidth(3.5));
label("$(5,12)$",(5,12),NE);label("$(12,5)$",(12,5),NE);label("$(13,0)$",(13,0),SE);label("$(0,13)$",(0,13)+(-0.9,0.3),NW);
draw((0,0)--(5,12),dashed);label("$13$",(2.5,6),E);
label("\mbox{}",(min(currentpicture,user=true)+max(currentpicture,user=true))/2,black); // keeps the text after the figure black when LaTeX typesets the labels inline
\end{asy}

(By symmetry only the first quadrant needs checking; the labelled points and their reflections in the axes give all twelve.)
